\section{Conclusions}\label{sec:concl}

We gave a complete account of resonant energy exchange through a gyroscopic
virtual mechanical system: design, failure and restoration.

\emph{Geometry.} Exact modal--polar coordinates reduce the exchange geometry to
energy, a log-amplitude ratio and a coprime relative phase.
In these coordinates the condition for complete host absorption is explicit,
and every resonant pair is absorbing at the right phase.
The mod-4 criterion of the conference version follows as the special case of
the impulse response.

\emph{Failure.} Any loop delay destabilises the interconnection, with a
closed-form count of unstable roots.
Sampling acts as an equivalent delay of one period.
Both pump the fast mode and drain the slow one, while static gain errors only
detune the phase.
Virtual-stiffness retuning tolerates only an $O(n^2)$ relative coupling error,
whereas a common-coupling trim has full authority.

\emph{Restoration.} Energy regulation, amplitude balancing with integral action
and phase locking through the common coupling give a triangular averaged system
that is ISS.
Its semiglobal practical ISS carries over to the actual loop in the
slow-adaptation regime.
Simulations confirm closed orbits and recurrent host depletion to $10^{-5}$ of the
initial energy under realistic imperfections.

The open items are listed in the gap markers.
The main ones are the rigorous inclusion of delay and sampling in the averaging
step, explicit gain bounds, the choice of the optimal absorbing target, and
experimental validation.
Extension to multi-degree-of-freedom hosts, where several host modes compete for
one virtual oscillator~\citep{juchem2024virtual}, is the natural next step.
